\section{«42»: citas filosóficas, destino y los límites del modelo del mundo}

El «42» del final viene de una broma de «Guía del viajero intergaláctico»: la respuesta a la pregunta fundamental de la vida, el universo y todo lo demás podría ser 42. Al cierre de la entrevista, el conductor pregunta «¿este mundo es un enorme modelo del mundo?» y «¿puedes predecir el destino?». Xie Saining (谢赛宁) responde que, por supuesto, el mundo puede verse como un enorme modelo del mundo, pero que no podemos predecir el destino porque no alcanzan los recursos; tal vez habría que usar la Tierra o el universo entero como computadora para obtener la respuesta.

Este pasaje no es un cierre en broma, sino un regreso a los límites del modelo del mundo: un world model existe para hacer predicciones suficientemente buenas con recursos limitados, no para copiar el universo de forma omnisciente y omnipotente. Cuanto más se acerca el modelo al mundo real, más tiene que enfrentar las limitaciones de recursos de cómputo, compresión de la representación, elección de objetivos e incognoscibilidad.

\subsection{Soltar a Wittgenstein}

Antes, el «42» sirvió para recordar que el modelo del mundo no puede fingir omnisciencia; aquí se vuelve a la relación entre el lenguaje y el mundo. Esta subsección importa porque muchos relatos sobre los LLM recurren justamente a frases filosóficas célebres para presentar la capacidad lingüística como comprensión del mundo.

La crítica de Xie Saining al uso de las citas de Wittgenstein es, en realidad, una forma de trazar los límites del modelo del mundo. La tesis del primer Wittgenstein sobre lenguaje y mundo tiene un contexto filosófico específico, y el Wittgenstein tardío se orientó después hacia los juegos de lenguaje. Tomar «los límites de mi lenguaje son los límites de mi mundo» como prueba directa de que los LLM pueden abarcar el mundo es un desajuste de contexto.

Si se sigue la idea tardía de los «juegos de lenguaje», según la cual el significado del lenguaje proviene de la práctica y de las formas de vida, esto se acerca más bien a la postura del modelo del mundo: el lenguaje no son tokens suspendidos en el aire, sino que adquiere significado al actuar, usarse y recibir retroalimentación en el mundo real.

\begin{importantbox}{La relación entre el lenguaje y el mundo}
El lenguaje puede comprimir el mundo, comunicarlo y organizar la experiencia, pero no puede sustituirlo. Lo que el modelo del mundo debe aprender es la parte de la estructura en la que el lenguaje se relaciona con la práctica: acción, retroalimentación, restricciones, causalidad y previsibilidad.
\end{importantbox}

\subsection{The normal one y la batería del equipo}

Por último, se vuelve a the normal one. Xie Saining retoma una expresión de Jürgen Klopp y se imagina como la batería del equipo: alguien que, con entusiasmo y energía, les da carga a los demás. No se trata de un discurso motivacional vacío, sino de una función central de una research organization. La exploración a largo plazo trae mucha frustración, fracasos e incertidumbre, y el equipo necesita a alguien que sostenga la dirección, la confianza y el ritmo.

También reconoce que el trasfondo de la investigación suele ser desalentador: los momentos de verdadera alegría quizá sean solo del 5\% al 10\%, cuando algo por fin sale. La afirmación es muy honesta y también explica por qué una línea como la del modelo del mundo necesita el respaldo de una cultura organizacional. Sin resiliencia psicológica a largo plazo, una visión ambiciosa pronto queda aplastada por los reveses de corto plazo.

\subsection{Resumen de la sección}

El cierre junta las preguntas técnicas y las preguntas de la vida: el mundo puede modelarse, pero el modelo siempre está limitado por los recursos y los objetivos; el lenguaje puede inspirar el pensamiento, pero no sustituye la práctica; una persona común también puede llegar a los grandes problemas mediante decisiones de largo plazo, la energía del equipo y el criterio de investigación.
